\documentclass[10pt,twocolumn]{article}
\usepackage[T1]{fontenc}
\usepackage[utf8]{inputenc}
\usepackage{lmodern}
\usepackage[margin=1.8cm,top=2.4cm,bottom=2cm,columnsep=0.6cm]{geometry}
\usepackage{fancyhdr}
\usepackage{titlesec}
\usepackage{booktabs}
\usepackage{amsmath,amssymb,amsfonts}
\usepackage{microtype}
\usepackage{xcolor}
\usepackage{mdframed}
\usepackage{parskip}
\usepackage{enumitem}
\usepackage{hyperref}

\definecolor{rulered}{RGB}{180,30,30}
\definecolor{boxgray}{RGB}{245,245,245}
\definecolor{keywordgray}{RGB}{100,100,100}

\pagestyle{fancy}
\fancyhf{}
\fancyhead[L]{\textsc{\small Claude Mag}}
\fancyhead[C]{\small\textit{Machine Learning Research Digest}}
\fancyhead[R]{\small 2026-07-15}
\fancyfoot[C]{\small\thepage}
\renewcommand{\headrulewidth}{0.8pt}

\titleformat{\section}{\normalsize\bfseries\color{rulered}}{}{0pt}{}%
  [\vspace{-4pt}\color{rulered}\rule{\columnwidth}{0.4pt}]
\titlespacing{\section}{0pt}{8pt}{4pt}

\hypersetup{colorlinks=true,urlcolor=rulered,linkcolor=black}

\begin{document}
\setlength{\parindent}{0pt}
\setlength{\parskip}{4pt}

{\small\color{keywordgray}\textbf{Keywords:}
  neural collapse \textbullet{}
  representation learning \textbullet{}
  information theory \textbullet{}
  language models}\par\vspace{4pt}

{\LARGE\bfseries\raggedright
  Neural Collapse Is Forbidden}\par\vspace{4pt}

{\large\itshape\raggedright
  An Information Floor Theorem Shows That Within-Class Variance
  in Language Models Is Obligatory, Not a Training Artifact}\par\vspace{6pt}

{\small\textbf{By} Bruno Abrahao
  (NYU Shanghai)}\par\vspace{2pt}

{\small\color{keywordgray}arXiv:2607.09487 \textbullet{}
  July 10, 2026}
\par\vspace{2pt}\noindent{\color{rulered}\rule{\columnwidth}{0.6pt}}\vspace{6pt}

Neural collapse---the phenomenon where final-layer representations
within each class converge to a single point during training---has been
hailed as evidence that deep networks find elegant, maximally-separated
class geometries. But for language models, this paper proves that
neural collapse is theoretically impossible: within-class variance
cannot shrink below a positive lower bound determined by how much
contextual information each token carries beyond its category label.
The paper terms this limit the \textbf{information floor}, derives it
analytically, and confirms it empirically across 14 models spanning
a 100$\times$ parameter range.

\begin{mdframed}[backgroundcolor=boxgray,linecolor=rulered,linewidth=1pt,
    innerleftmargin=6pt,innerrightmargin=6pt,
    innertopmargin=6pt,innerbottommargin=6pt]
\textbf{\small AT A GLANCE}\par\vspace{4pt}
\begin{description}[leftmargin=0pt,labelsep=4pt,itemsep=2pt,parsep=0pt,
    font=\small\bfseries]
  \item[Problem:] Within-class variance in LM representations is
    commonly attributed to incomplete training or suboptimal
    regularization.
  \item[Why it matters:] If within-class variance is obligatory, attempts
    to reduce it via regularization or architecture choices will hurt
    model performance.
  \item[Method:] Probing study across 14 models; theoretical proof of
    an information floor on within-class variance; k-NN mutual information
    estimator to compute the bound empirically.
  \item[Result:] Within-class variance accounts for 79--91\% of total
    representational variance; the theoretical floor covers 71--79\%.
  \item[Evidence:] 26-page paper with 3 figures and 19 tables across
    five text domains and five model families.
\end{description}
\end{mdframed}

\section{The Big Picture}

Neural collapse was first described in classification networks: as
training progresses, the penultimate-layer representations of samples
within each class collapse to a single class mean, and the class means
arrange themselves in a maximally-equidistant ``simplex ETF'' geometry.
The phenomenon has been viewed as a signature of an ideal training
outcome---a fully collapsed network is maximally linear-separable and
robust to perturbations.

For language models, the analogue of a ``class'' is ambiguous. Prior
work has used token categories (part of speech, semantic class, domain
label) as the class variable and found that within-category variance
is large---sometimes 80--90\% of total variance---and interpreted this
as evidence that language models fail to achieve neural collapse.

This paper argues the interpretation is exactly backwards: large
within-class variance is not a failure but a necessity. Language models
must encode token-specific contextual information within each category
because two tokens of the same category (e.g., both nouns) carry
different next-token distributions depending on their context.
Collapsing representations within a category would destroy this
information and increase prediction loss---so the training objective
itself prevents neural collapse.

\section{Why Existing Approaches Fall Short}

The neural collapse literature for classification assumes a fixed
label space where two samples of the same class carry identical
distributional targets. This assumption fails for language models:
the distribution over next tokens for ``bank'' depends on whether
the preceding context is financial or geographic, even though ``bank''
is in the same syntactic/semantic category in both cases.

Attempts to fine-tune language models toward neural collapse (via
additional regularization or loss terms) have been tried in the
context of classification fine-tuning. The paper's theoretical result
explains why these attempts always succeed in reducing within-class
variance on the fine-tuning data but increase perplexity on
held-out data: they destroy the contextual information encoded in
the within-class variance.

\section{Core Mechanism}

The information floor arises from the structure of next-token
prediction. For a token $X$ with category label $C$ and context $Z$,
the model must compute $p(X_{\text{next}} \mid X, Z)$, which depends
on both the category and the specific context. Two tokens of the same
category in different contexts have different conditional next-token
distributions, and these differences must be encoded in the
representation.

The conditional mutual information $I(X; Z \mid C)$ quantifies exactly
how much contextual information is necessary beyond the category label.
The paper proves that the within-class variance $\sigma^2_W$ in any
representation that achieves training loss below the theoretical
minimum must satisfy $\sigma^2_W \geq c \cdot I(X; Z \mid C)$.

\section{Technical Formulation}

Let $X$ be the token, $C$ its macro-category, $Z$ the context window,
and $h(X, Z) \in \mathbb{R}^d$ the final-layer representation. Define
the within-class variance as:
\[
\sigma^2_W = \mathbb{E}_C \bigl[
  \mathrm{Var}(h(X, Z) \mid C)
\bigr]
\]

\textbf{Theorem (Information Floor):}
For any model that achieves cross-entropy loss
$\mathcal{L} \leq H(X_{\text{next}} \mid C, Z)$ on the training
distribution:
\[
\sigma^2_W \;\geq\; c \cdot I(X;\, Z \mid C)
\]
where $c > 0$ depends only on the loss function curvature and
projection constants of the final linear layer.

The proof proceeds by lower-bounding the Fisher information of the
output distribution as a function of $\sigma^2_W$, and showing that
$\sigma^2_W < c \cdot I(X; Z \mid C)$ implies the Fisher information
is insufficient to support the observed training loss.

\section{Learning or Inference Procedure}

\begin{enumerate}[itemsep=2pt,parsep=0pt]
  \item For each of the 14 models, collect final-layer representations
    on 50,000 tokens from five text domains.
  \item Assign 32-category macro-category labels using WordNet +
    POS tagging.
  \item Compute within-class and between-class variance in the
    top-$k$ principal subspace of the full covariance matrix.
  \item Estimate $I(X; Z \mid C)$ using a k-NN mutual information
    estimator (Kozachenko-Leonenko calibrated).
  \item Compare observed $\sigma^2_W / (\sigma^2_B + \sigma^2_W)$
    against the theoretical lower bound derived from $I(X; Z \mid C)$.
\end{enumerate}

\section{What the Guarantee Says}

The information floor theorem provides a constructive lower bound:
given an estimate of $I(X; Z \mid C)$ from data, the bound predicts
the minimum observable within-class variance share. Across five
domains and 14 models, the theoretical floor covers 71--79\% of
total variance, and the observed share is 79--91\%---consistently
above the floor by a margin of 8--12 percentage points, consistent
with the floor being a lower bound rather than an equality.

\section{Experimental Findings}

\begin{center}
\small
\begin{tabular}{lcc}
\toprule
Model family & $\sigma^2_W / \sigma^2_{\text{total}}$ & Floor \\
\midrule
GPT series & $82.4\% \pm 3.1\%$ & 73.1\% \\
LLaMA series & $84.7\% \pm 2.4\%$ & 75.8\% \\
Mistral & $83.9\% \pm 2.8\%$ & 74.2\% \\
Claude family & $79.3\% \pm 4.2\%$ & 71.0\% \\
Gemma & $85.2\% \pm 1.9\%$ & 76.3\% \\
\bottomrule
\end{tabular}
\end{center}

The variance share is insensitive to scale within each family
(max 2\% change over 10$\times$ parameter increase) and insensitive
to tokenization (BPE vs.\ SentencePiece: 0.8\% difference).
All observations lie above the predicted floor.

\section{Ablations and Interpretation}

\textbf{Random-label control:} Replacing token category labels with
random assignments increases $\sigma^2_B$ by 18\% and $\sigma^2_W$
by 22\%, confirming the variance decomposition is meaningful.

\textbf{Layer depth:} Within-class variance share increases from 55\%
at layer 1 to 84\% at the final layer, suggesting the model
progressively specializes representations toward context-specific
rather than category-level structure.

\textbf{Fine-tuning toward collapse:} Fine-tuning on sentiment
classification with an additional within-class variance penalty reduces
$\sigma^2_W$ on the fine-tuning domain by 12\% but increases
perplexity by 8\%, directly confirming that forcing collapse destroys
necessary contextual information.

\textbf{Rare tokens:} The information floor is higher for rare tokens
(those with unique distributional contexts), consistent with the
finding that rare tokens are disproportionately memorized---they
carry more $I(X; Z \mid C)$ and therefore require more within-class
dispersion to represent correctly.

\section*{Reference}

Bruno Abrahao.
\textbf{Neural Collapse Is Forbidden: Information Floors in
Language Models}.
arXiv:2607.09487, July 10, 2026. NYU Shanghai.
\url{https://arxiv.org/abs/2607.09487}

\end{document}
